\chapter{Constitutive models}
\label{ch:constitutive}
\impl{src/physics.jl, src/materials.jl}

Carina delegates constitutive evaluation to \code{ConstitutiveModels.jl}.  The
interface is narrow by design: a model returns $\Piola$ given $\Grad\disp$ and,
for models carrying internal variables, an updated state.  Four models are
wired into the input parser.

\section{Linear elastic}

Small-strain isotropic elasticity, $\Cauchy = \lambda\tr(\strain)\Ident +
2\mu\strain$ with $\strain = \tfrac12(\Grad\disp +
\Grad\disp^{\mathsf{T}})$.  Its tangent is constant, which Carina exploits: for
a linear-elastic model the stiffness is assembled once and the factorization
--- or, on GPU, the Cholesky factor uploaded for triangular solves
(Section~\ref{sec:gpu-cholesky}) --- is cached for the entire run.

\section{Neo-Hookean}

The compressible neo-Hookean strain energy
\begin{equation}
  \Psi(\Fdef) = \frac{\mu}{2}\left(\tr\Cright - 3\right) - \mu \ln J
              + \frac{\lambda}{2}\left(\ln J\right)^2
\end{equation}
gives
\begin{equation}
  \Piola = \mu\left(\Fdef - \Fdef^{-\mathsf{T}}\right)
         + \lambda \ln J \, \Fdef^{-\mathsf{T}} .
\end{equation}
This is the workhorse for the benchmark problems throughout the performance
campaigns, and the one model whose tangent \emph{action} Carina writes in
closed form (Section~\ref{sec:tangent-strategies}).

\section{Hencky}

Logarithmic (true) strain elasticity, $\Kirch = \lambda\tr(\vect{h})\Ident +
2\mu\vect{h}$ with $\vect{h} = \tfrac12\ln\Cright$.  The matrix logarithm is
evaluated by inverse scaling-and-squaring with Pad\'e approximants rather than
by spectral decomposition: the spectral route loses accuracy and
differentiability when eigenvalues of $\Cright$ coalesce, which is exactly the
neighbourhood of the undeformed state.

\section{Finite-deformation \texorpdfstring{$J_2$}{J2} plasticity}

Rate-independent $J_2$ flow with linear isotropic hardening, integrated by a
radial-return map at each quadrature point.  The trial elastic state is
computed, and if the yield function
\begin{equation}
  f = \norm{\dev \Kirch^{\mathrm{tr}}} - \sqrt{\tfrac23}\left(\sigma_y + H\,\bar{\varepsilon}^p\right)
\end{equation}
is positive the plastic multiplier is solved for and the stress projected back
onto the yield surface.  The model carries internal state --- the plastic
strain tensor and the equivalent plastic strain --- which makes it
\emph{path dependent}: the same displacement reached by different increments
gives different stresses.  This has two consequences that surface elsewhere in
this manual.  Large load steps miss the yield surface and must be prevented by
adaptive stepping (Section~\ref{sec:adaptive}); and the return map can raise
domain errors on non-physical trial states, which the solver treats as a step
failure rather than a crash.

\section{Tangent strategies}
\label{sec:tangent-strategies}

Equation~\eqref{eq:tangent} needs $\fourth{A} = \partial\Piola/\partial
\Grad\disp$ at every quadrature point.  Carina obtains it three different ways,
chosen by what the model can support.

\paragraph{Closed form.}  For neo-Hookean the directional derivative is
written analytically.  With $\delta\Fdef = \Grad\vect{v}$, and abbreviating
$a = \tr(\Fdef^{-1}\delta\Fdef)$ and $\tensor{W} =
\Fdef^{-\mathsf{T}}\delta\Fdef^{\mathsf{T}}\Fdef^{-\mathsf{T}}$,
\begin{equation}
  \delta\Piola = \mu\left(\delta\Fdef + \tensor{W}\right)
               + \lambda\, a\, \Fdef^{-\mathsf{T}}
               - \lambda \ln J\, \tensor{W} .
  \label{eq:nh-action}
\end{equation}
This form is what the matrix-free operator of Chapter~\ref{ch:matrixfree}
evaluates: it needs the tangent's action on \emph{one} direction, never the
full $81$-component $\fourth{A}$, so the closed form costs a handful of
$3\times3$ products instead of a rank-four object.

\paragraph{Automatic differentiation.}  For stateless hyperelastic models
generally, $\fourth{A}$ is obtained by differentiating \code{pk1\_stress}
through dual numbers.  This is exact to machine precision and requires nothing
from the model author.

\paragraph{Finite differences.}  Models carrying state --- $J_2$ plasticity ---
run a return map inside \code{pk1\_stress}, whose branches are not
differentiable in the ordinary sense.  Carina falls back to a forward finite
difference over the nine components of $\Grad\disp$,
\begin{equation}
  \fourth{A}_{ijkl} \approx
  \frac{\Piola_{ij}(\Grad\disp + h_{kl}\,\vect{e}_k\otimes\vect{e}_l) - \Piola_{ij}(\Grad\disp)}{h_{kl}},
\end{equation}
at a cost of nine stress evaluations per quadrature point.  The perturbation is
adaptive,
\begin{equation}
  h_{kl} = \max\left(\sqrt{\varepsilon_{\mathrm{mach}}}\,\abs{(\Grad\disp)_{kl}},\ h_{\min}\right),
  \qquad h_{\min} = 10^{-10},
\end{equation}
following the Sierra/SM convention.  The relative term dominates wherever the
strain is resolved; the floor prevents a vanishing perturbation --- and hence
catastrophic cancellation --- at components that happen to be near zero, which
in an undeformed region is most of them.

\begin{remark}
The finite-difference tangent is a \emph{tangent}, not the residual: it enters
only the Newton direction and the preconditioner, never the convergence test.
An approximate tangent slows Newton's asymptotic rate but cannot change the
solution it converges to, because that is fixed by $\Rvec = \vect{0}$.  The
same argument licenses the reduced-precision operator of
Section~\ref{sec:mixed-precision}.
\end{remark}
